\scanpage{447}
\mspage{359}
qui est $L$-$\pi$-adaptée (cf \unsure{page|pages} 343, c), on a ceci : $x\in\mathrm{Ob}\,S$ sss $I_{\underline S}$ est cofinal dans $I$.

c$'$) Comme c, avec $I$ \unsure{ordonnée|ordonné}.

\inserted{[ c$''$) (Si on a choisi, pour tt $x\in\mathrm{Ob}\,M$, \uline{une} représentation $\pi$-adaptée $x\simeq\underset{I}{\underrightarrow{\lim}}\,x_i$) : $x\in\mathrm{ob}\,\underline S\Longleftrightarrow$ \unsure{l'ens|alors} $I_{\underline S}$ est cofinal dans $I$ ]}

\margin{\deleted{\unsure{B}}}
\uline{Démonstration} C'est essentiellement l'énoncé de 4.13.22 dans le cas particulier où $\Phi={}$ ens des catégories \uline{filtrantes} de cardinal $\leq\pi$, mais où, pour simplifier, on a supposé d'avance que $\underline S(\pi)$ \unsure{satisfasse|satisfais|satisfait} à $L(\Phi)$. L'équivalence de a) b) c) est contenue dans \unsure{4.13.??|4.13.22}, ainsi que $c\Longrightarrow c'$. Mais \deleted{il faut} pour $c'\Longrightarrow a$, \deleted{\unsure{les|le}} l'hypothèse \inserted{accessoire} $\pi^{\unsure{\bullet|\circ|'}}=\pi$ et \deleted{le fait} que les $J\in\Phi$ sont grdes devant $\pi_0$, n'est pas satisfaite ici, et il faut une démonstration indépendante, calquée sur celle des \deleted{\unsure{les}} pages 333-336. \deleted{Dans \unsure{ce} résultat cette démonstration} Il suffit pour cel\unsure{a}, à l'endroit $\square$ page 333, d'invoquer le

\uline{Lemme} 4.13.37. \unsure{Soient|Soit} $M$ une \unsure{cat-|cat.}
% archive sheet number: 446
